\chapter*{Z}
\label{chap:z}

\term{Z-score}
A measure of how many standard deviations a data point is from the mean of a distribution: $z = (x - \mu) / \sigma$. It is used to standardize different datasets for comparison and is the basis for the z-test in statistics.

\term{Z-test}
A statistical hypothesis test used to determine whether the means of two populations are different when the population variance is known and the sample size is large.

\term{Z-transform}
A discrete-time version of the Laplace transform, mapping a sequence $\{x[n]\}$ to a complex function $X(z) = \sum_{n=-\infty}^\infty x[n] z^{-n}$. It is widely used in digital signal processing and control theory to analyze linear time-invariant systems.

\term{Zeros of a Function}
The set of points $z$ in the domain of a function $f$ such that $f(z) = 0$. In complex analysis, zeros are classified as isolated or non-isolated, and their order (multiplicity) is a key feature of the function's behavior.

\term{Zeta Function (Riemann)}
The function $\zeta(s) = \sum_{n=1}^\infty n^{-s}$ for $\operatorname{Re}(s) > 1$, analytically continued to the whole complex plane. The Riemann Hypothesis, stating that all non-trivial zeros lie on the critical line $\operatorname{Re}(s) = 1/2$, is one of the most important open problems in mathematics.

\term{Zigzag}
A polygonal path alternating direction sharply, used in graph theory and combinatorics as a pattern of signs or steps.

\term{Zonal Polynomial}
A symmetric homogeneous polynomial appearing in multivariate statistics and representation theory, associated with spherical functions.

\term{Zonohedron}
A convex polyhedron whose faces are centrally symmetric, a Minkowski sum of line segments.

\term{Zorn's Lemma}
A statement in set theory: every partially ordered set in which every chain has an upper bound contains at least one maximal element. It is equivalent to the Axiom of Choice and is used to prove the existence of bases in every vector space and maximal ideals in rings.

\term{Zygmund Class}
A function belongs to a Zygmund class when its second difference satisfies a bound such as $|f(x+h)+f(x-h)-2f(x)|\le C|h|$ for all relevant $x$ and $h$. This describes a form of smoothness close to, but slightly weaker than, Lipschitz continuity. Applications: Zygmund classes help analyze Fourier series, approximation error, and the behavior of functions near singularities.

\term{Zeno's Paradoxes}
A set of philosophical problems, attributed to Zeno of Elea, illustrating the difficulties of conceptualizing infinity, motion, and continuity. The Dichotomy paradox suggests an object cannot reach a destination because it must first reach the halfway point, then the halfway point of the remainder, and so on. Achilles and the Tortoise argues that the swift runner Achilles never overtakes a tortoise given a head start, since each gain leaves a further gap. The Arrow paradox claims a flying arrow occupies a fixed position at every instant, so motion is impossible. The Stadium paradox concerns the relative motion of rows of bodies in the smallest indivisible instant. The first two are resolved by the convergence of infinite geometric series; the last two sharpen the distinction between continuous motion and its discrete mathematical description.

\term{Zappa--Szép Product}
A generalization of the semi-direct product of groups where two groups $G$ and $H$ act on each other. The resulting group $G \bowtie H$ is the set $G \times H$ with a multiplication rule that incorporates both actions.

\term{Zariski Topology}
A topology used in algebraic geometry where the closed sets are the zero-sets of collections of polynomials. In this topology, open sets are very large (dense), making the space non-Hausdorff.

\term{Zariski Closure}
The smallest algebraic variety (the set of zeros of a set of polynomials) containing a given subset of an affine or projective space. It is the closure of the set in the Zariski topology.

\term{Z-module}
A module over the ring of integers $\mathbb{Z}$. Every abelian group is naturally a $\mathbb{Z}$-module, making the study of $\mathbb{Z}$-modules identical to the study of abelian groups.

\term{Zero Divisor}
An element $a \neq 0$ in a ring $R$ such that there exists another non-zero element $b \in R$ with $ab = 0$. An integral domain is defined as a commutative ring with no zero divisors.

\term{Zero-sum Game}
A game in game theory where the sum of the gains and losses of all players is zero. In such games, one player's gain is exactly balanced by the other players' losses.

\term{Zero-one Law}
A phenomenon in mathematical logic and probability where the probability of a first-order sentence being true in a random structure is either 0 or 1 as the size of the structure tends to infinity.

\term{Zero-set}
The set of all points in the domain of a function where the function's value is zero. In topology, zero-sets are used to define the fact that a space is completely regular.

\term{Zero-point Energy}
The lowest possible energy that a quantum mechanical physical system may have. In the context of the harmonic oscillator, this is $E_0 = \frac{1}{2}\hbar\omega$.

\term{Zygmund's Theorem}
A result in harmonic analysis concerning the convergence of Fourier series of functions in the Zygmund class, showing that such functions possess a higher degree of regularity than general continuous functions.

\end{document}
